\chapter{Sesión 3: Seguridad JWT y autenticación}
\label{ch:sesion3}

\section{Objetivo de la sesión}

Implementar el sistema de autenticación basado en JSON Web Tokens (JWT), cubriendo:

\begin{itemize}[leftmargin=*]
  \item Registro de usuarios (\texttt{POST /api/auth/register}).
  \item Inicio de sesión (\texttt{POST /api/auth/login}).
  \item Renovación de token (\texttt{POST /api/auth/refresh}).
  \item Filtro que valida el JWT en cada petición autenticada.
  \item Configuración de Spring Security en modo \textit{stateless}.
  \item Hash de contraseñas con BCrypt.
\end{itemize}

\section{Contexto}

En la Sesión~\ref{ch:sesion2} el backend arrancaba pero Spring Security generaba una contraseña temporal y todos los endpoints (salvo los públicos por defecto) quedaban bloqueados con \texttt{inMemoryUserDetailsManager}. No existía forma de crear usuarios en la base de datos ni de autenticarse contra ella.

\section{Desarrollo}

\subsection{Estructura creada}

\begin{lstlisting}[language=bash]
src/main/java/mx/uam/notiuam/
|- config/
|  `- SecurityConfig.java
|- security/
|  |- JwtService.java
|  |- JwtAuthenticationFilter.java
|  `- UserDetailsServiceImpl.java
|- dto/auth/
|  |- RegisterRequest.java
|  |- LoginRequest.java
|  |- RefreshRequest.java
|  `- AuthResponse.java
|- service/
|  `- AuthService.java
|- controller/
|  `- AuthController.java
`- exception/
   `- GlobalExceptionHandler.java
\end{lstlisting}

\subsection{Decisiones técnicas}

\begin{itemize}[leftmargin=*]
  \item \textbf{JWT stateless.} No se usan sesiones HTTP. Toda la identidad viaja en el token firmado con HMAC-SHA256.
  \item \textbf{Dos tipos de token.} \textit{Access token} de 1 hora y \textit{refresh token} de 7 días. El claim \texttt{type} distingue uno del otro.
  \item \textbf{BCrypt.} Las contraseñas se almacenan siempre hasheadas, nunca en claro.
  \item \textbf{Alias \texttt{ROLE\_}.} Los roles se exponen como \texttt{ROLE\_STUDENT}, \texttt{ROLE\_PROFESSOR}, \texttt{ROLE\_ADMIN} para que \texttt{@PreAuthorize} funcione sin configuración extra.
  \item \textbf{CORS abierto en desarrollo.} Se permite cualquier origen (\texttt{*}) para que la app Flutter pueda probar sin configuración adicional. En producción se restringirá.
\end{itemize}

\subsection{Componentes clave}

\subsubsection{\texttt{JwtService}}

Genera, firma y valida tokens. Métodos principales:

\begin{lstlisting}[language=java]
public String generateAccessToken(User user);
public String generateRefreshToken(User user);
public Claims parse(String token);
public boolean isValid(String token);
public long getAccessExpirationMs();
\end{lstlisting}

El claim \texttt{subject} guarda el \texttt{UUID} del usuario. Otros claims: \texttt{email}, \texttt{role}, \texttt{type} (\texttt{access} o \texttt{refresh}).

\subsubsection{\texttt{JwtAuthenticationFilter}}

Extiende \texttt{OncePerRequestFilter}. En cada petición:
\begin{enumerate}[leftmargin=*]
  \item Lee el header \texttt{Authorization: Bearer <token>}.
  \item Valida el token con \texttt{JwtService}.
  \item Extrae \texttt{userId}, \texttt{email}, \texttt{role}.
  \item Coloca un \texttt{UsernamePasswordAuthenticationToken} en el \texttt{SecurityContext}.
  \item El principal es un record \texttt{AuthenticatedUser(UUID id, String email, String role)}.
\end{enumerate}

\begin{nota}[Nota]
Si el token no es válido o falta, el filtro simplemente deja pasar la petición sin autenticación. La decisión de permitirla o rechazarla la toma la cadena de filtros de Spring Security según el matcher de la URL.
\end{nota}

\subsubsection{\texttt{SecurityConfig}}

Define:
\begin{itemize}[leftmargin=*]
  \item \texttt{PasswordEncoder} con BCrypt.
  \item \texttt{AuthenticationManager} respaldado por \texttt{DaoAuthenticationProvider} + \texttt{UserDetailsServiceImpl}.
  \item \texttt{SecurityFilterChain}:
    \begin{itemize}
      \item CSRF deshabilitado (API stateless).
      \item Sesiones en \texttt{STATELESS}.
      \item \texttt{/api/auth/**}, Swagger y \texttt{/actuator/health} públicos.
      \item Todo lo demás, autenticado.
      \item \texttt{JwtAuthenticationFilter} antes de \texttt{UsernamePasswordAuthenticationFilter}.
    \end{itemize}
\end{itemize}

\subsubsection{\texttt{AuthService}}

Contiene la lógica de negocio:

\begin{itemize}[leftmargin=*]
  \item \texttt{register}: valida unicidad del email, hashea la contraseña, guarda y devuelve par de tokens.
  \item \texttt{login}: usa el \texttt{AuthenticationManager} para validar credenciales, luego emite tokens.
  \item \texttt{refresh}: valida el refresh token, verifica el claim \texttt{type=refresh} y emite un nuevo par.
\end{itemize}

\subsection{\texttt{AuthController}}

Expone los tres endpoints bajo \texttt{/api/auth}:

\begin{lstlisting}[language=java]
@PostMapping("/register")  // 201 Created
@PostMapping("/login")     // 200 OK
@PostMapping("/refresh")   // 200 OK
\end{lstlisting}

Todas las peticiones usan \texttt{@Valid} para disparar la validación de los DTOs.

\subsection{Manejo global de errores}

\texttt{GlobalExceptionHandler} con \texttt{@RestControllerAdvice} convierte:

\begin{itemize}[leftmargin=*]
  \item \texttt{IllegalArgumentException} $\rightarrow$ \texttt{400 Bad Request}.
  \item \texttt{MethodArgumentNotValidException} $\rightarrow$ \texttt{400} con detalle de campos inválidos.
\end{itemize}

Respuesta uniforme:

\begin{lstlisting}
{
  "timestamp": "2026-10-05T22:30:00Z",
  "status": 400,
  "error": "Bad Request",
  "message": "El correo ya esta registrado"
}
\end{lstlisting}

\subsection{Ajustes al \texttt{application.yml}}

Se limpiaron dos warnings que aparecieron en la Sesión~\ref{ch:sesion2}:

\begin{itemize}[leftmargin=*]
  \item Se eliminó \texttt{hibernate.dialect} (Hibernate 7 lo autodetecta).
  \item Se añadió \texttt{spring.data.redis.repositories.enabled: false} para evitar el escaneo de repositorios Redis.
\end{itemize}

\section{Verificación}

\subsection{Registro}

\begin{lstlisting}[language=bash]
curl -s -X POST http://localhost:8081/api/auth/register \
  -H "Content-Type: application/json" \
  -d '{
    "email": "aaron@uam.mx",
    "password": "password123",
    "displayName": "Aaron Ramos",
    "role": "STUDENT"
  }' | jq
\end{lstlisting}

\begin{exito}[Respuesta esperada]
Un JSON con \texttt{accessToken}, \texttt{refreshToken}, \texttt{tokenType}, \texttt{expiresInMs} y un objeto \texttt{user} con los datos del usuario creado.
\end{exito}

\subsection{Login}

\begin{lstlisting}[language=bash]
curl -s -X POST http://localhost:8081/api/auth/login \
  -H "Content-Type: application/json" \
  -d '{"email":"aaron@uam.mx","password":"password123"}' | jq
\end{lstlisting}

Debe devolver un par de tokens nuevos.

\subsection{Refresh}

\begin{lstlisting}[language=bash]
REFRESH=$(curl -s -X POST http://localhost:8081/api/auth/login \
  -H "Content-Type: application/json" \
  -d '{"email":"aaron@uam.mx","password":"password123"}' | jq -r .refreshToken)

curl -s -X POST http://localhost:8081/api/auth/refresh \
  -H "Content-Type: application/json" \
  -d "{\"refreshToken\":\"$REFRESH\"}" | jq
\end{lstlisting}

Debe emitir un nuevo par de tokens.

\subsection{Verificar hash en PostgreSQL}

\begin{lstlisting}[language=bash]
docker exec -it notiuam-postgres psql -U notiuam -d notiuam \
  -c "SELECT email, display_name, role,
             LEFT(password_hash, 7) AS hash_prefix
      FROM users;"
\end{lstlisting}

\begin{exito}[Resultado esperado]
El prefijo del hash debe ser \texttt{\$2a\$10\$} (BCrypt). Nunca aparece la contraseña en claro.
\end{exito}

\subsection{Comprobar que arranca sin contraseña temporal}

Al reiniciar el backend ya \textbf{no} debe aparecer el mensaje:

\begin{lstlisting}[language=bash]
Using generated security password: ...
\end{lstlisting}

Si aparece, revisar que \texttt{UserDetailsServiceImpl} esté en el paquete \texttt{mx.uam.notiuam.security} y que la dependencia \texttt{spring-boot-starter-security} esté en el \texttt{pom.xml}.

\section{Resultado}

\begin{exito}[Estado al cierre]
\begin{itemize}[leftmargin=*]
  \item Registro, login y refresh funcionales.
  \item Autenticación stateless con JWT (access 1h + refresh 7d).
  \item Contraseñas hasheadas con BCrypt.
  \item Roles expuestos como \texttt{ROLE\_*} para usar \texttt{@PreAuthorize}.
  \item CORS listo para el cliente Flutter.
  \item Manejo global de errores unificado.
\end{itemize}
\end{exito}

\textbf{Pendiente para la siguiente sesión:} endpoints de canales, suscripciones y publicaciones, usando el \texttt{AuthenticatedUser} del \texttt{SecurityContext} para identificar al emisor.